\section{Plan de pruebas y experimentaci\'on}
\label{sec:pruebas}

Este plan define c\'omo se demostrar\'a la correcci\'on funcional, la
seguridad frente a manipulaciones y el desempe\~no del sistema; la
implementaci\'on y las mediciones se presentar\'an en el Entregable~2. Todas
las pruebas del Cuadro~\ref{tab:plan-pruebas} ser\'an automatizadas.

\begin{table}[ht]
\centering
\small
\begin{tabularx}{\textwidth}{@{}l p{6.2cm} X@{}}
\toprule
Tipo & Acci\'on de prueba & Comportamiento esperado \\
\midrule
Casos v\'alidos &
Pedir el primer registro ($i=1$), el \'ultimo ($i=N$) y varios al azar;
repetir el mismo $i$ con una sesi\'on nueva. &
El registro descifrado coincide byte a byte con el original; $sid$, $u$, $v$
y $\{E_j\}$ son distintos en cada ejecuci\'on. \\
Claves no pedidas &
Intentar abrir $E_{j\ne i}$ con la clave derivada de $Z_i$. &
AEAD rechaza con \code{InvalidTag}. \\
Entradas inv\'alidas &
Enviar un $u$ no can\'onico; reutilizar un $sid$ ya consumido; usar un $sid$
vencido; pedir una versi\'on inexistente; pedir un \code{record\_id}
inexistente o un \'indice fuera de $[1,N]$. &
HTTP 400, 409, 404 y 422, en ese orden. El \'indice inv\'alido lo rechaza el
cliente sin enviar nada, porque el servidor nunca ve $i$. \\
Manipulaci\'on &
Cambiar un bit en un registro $C_j$ del cat\'alogo, en una encapsulaci\'on
$E_j$, en un nonce o en los datos asociados. &
\code{sha256\_root} detecta el cambio en el cat\'alogo y el tag AEAD rechaza
los dem\'as (\code{InvalidTag}). \\
Ataque de desfase &
Un proxy entre cliente y servidor sustituye $u$ por $\hat u = u \cdot v$
(Ej.~11.20); en una variante, adem\'as desplaza una posici\'on la lista de
encapsulaciones. &
El cliente no acepta ning\'un registro: abrir $E_i$ falla con
\code{InvalidTag} en ambas variantes. \\
\bottomrule
\end{tabularx}
\caption{Plan de pruebas funcionales, negativas y de manipulaci\'on.}
\label{tab:plan-pruebas}
\end{table}

\subsection{Metodolog\'ia de experimentaci\'on}
Para medir el costo de ocultar la selecci\'on se evaluar\'an cat\'alogos de
$N \in \{500, 1000, 2500, 5000\}$ registros de 256\,B, comparando el modo OT
con la l\'inea base HTTP en la misma m\'aquina: una estaci\'on x86\_64
(Intel Core i7 o AMD Ryzen con AES-NI, 16\,GB de RAM) con Python~3.11. Los
tiempos se toman con \code{time.perf\_counter\_ns()}, con 5 repeticiones por
configuraci\'on para reportar media y desviaci\'on est\'andar.

Un script (\code{bench.py}) registrar\'a en CSV: (1) latencia de extremo a
extremo; (2) tiempo de c\'omputo en cliente y servidor; (3) operaciones de
grupo (cliente: 3 exponenciaciones y 1 multiplicaci\'on; servidor: $N+1$ y
$N$); y (4) bytes de la petici\'on, la respuesta y la descarga inicial. El
Cuadro~\ref{tab:plantilla-experimentos} adelanta los tama\~nos esperados
seg\'un los formatos de la Secci\'on~\ref{sec:formatos}; el Entregable~2
los contrastar\'a con lo medido.

\begin{table}[ht]
\centering
\small
\begin{tabularx}{\textwidth}{@{}r X X X X X@{}}
\toprule
$N$ & Descarga inicial & Petici\'on OT & Respuesta OT & Respuesta normal & Latencia (ms) \\
\midrule
500  & $\approx 0.22$\,MB & 96\,B & $\approx 55$\,KB  & $\approx 0.3$\,KB & \emph{A medir en E2} \\
1000 & $\approx 0.44$\,MB & 96\,B & $\approx 111$\,KB & $\approx 0.3$\,KB & \emph{A medir en E2} \\
2500 & $\approx 1.10$\,MB & 96\,B & $\approx 279$\,KB & $\approx 0.3$\,KB & \emph{A medir en E2} \\
5000 & $\approx 2.19$\,MB & 96\,B & $\approx 559$\,KB & $\approx 0.3$\,KB & \emph{A medir en E2} \\
\bottomrule
\end{tabularx}
\caption{Tama\~nos estimados de los cuerpos JSON (sin cabeceras HTTP) y plantilla para las latencias OT y normal del Entregable~2.}
\label{tab:plantilla-experimentos}
\end{table}
